\section{Implementation Details}

\subsection{Synthetic cases and numerical solver}

All reported experiments use the 1D periodic solver in the accompanying codebase. Trajectories are generated by an RK4 time integrator applied to a conservative finite-volume-style discretization of
\begin{equation}
u_t = \partial_x \left( D(u) u_x \right) + R(u).
\end{equation}
The three main benchmark cases are
\begin{align}
\text{Case A:}\quad & D(u) = 0.01 + 0.05u, \qquad R(u) = u(1-u), \\
\text{Case B:}\quad & D(u) = 0.01 + 0.03u^2, \qquad R(u) = u - 1.5u^2 + 0.5u^3, \\
\text{Case Exp:}\quad & D(u) = 0.01 + 0.035(1-e^{-2.5u}), \qquad R(u) = u(1.1e^{-1.4u} - 0.22).
\end{align}
Case~A and Case~B are matched to the low-order polynomial baseline library. Case~Exp is a smooth non-polynomial stress test for function-class mismatch.

\subsection{Training and observation protocol}

Unless otherwise stated, the baseline and symbolic benchmarks use $8$ training trajectories, random smooth Fourier initial conditions with $4$ modes, and the simulation grid
\begin{equation}
n_x = 64,\qquad \Delta t = 10^{-4},\qquad T = 0.1,
\end{equation}
with snapshots saved every $10$ time steps. The excitation comparison in Table~\ref{tab:excitation} uses a lighter protocol with $n_x=48$, $T=0.05$, snapshots saved every $5$ steps, and $6$ training trajectories. The cross-resolution benchmark in Table~\ref{tab:cross_resolution} uses the explicit fine and validation grids reported in its caption. The main surrogate is a two-branch MLP with hidden width $64$, depth $2$, and $150$ epochs of Adam optimization at learning rate $10^{-3}$.

Observation degradation is applied after simulation. A ``5\% noise'' setting means additive Gaussian noise with standard deviation $0.05\,\mathrm{std}(u)$, followed by clipping to the admissible state range. A ``sparse $x+t$'' setting means spatial stride $2$ and temporal stride $2$ in the observed trajectories. All tables report mean$\pm$standard deviation over three random seeds unless stated otherwise.

\subsection{Restricted symbolic families}

The symbolic stage is intentionally restricted. For diffusion, the candidate library contains polynomial families up to degree $3$, a $(2,2)$ rational family, an exponential-decay family, and a saturation family. For reaction, the candidate library contains polynomial families up to degree $4$, a $(2,2)$ rational family, a $u\exp(\cdot)$ family, and a $u/(1+\cdot)$ saturation family. The selected expression is the lowest-complexity candidate whose fit mean-squared error lies within a $5\%$ relative tolerance and a $10^{-8}$ absolute tolerance of the best candidate. This library is not meant to be exhaustive; it is a controlled interpretable compression family.

\subsection{Reproducibility notes}

The main result tables are generated from CSV summaries written by the benchmark scripts in the accompanying codebase, including symbolic-compression, excitation, and cross-resolution runs. Fixed random seeds are used throughout. Exact bitwise determinism across hardware is not claimed because PyTorch deterministic kernels are not globally forced, but repeated runs in the current environment reproduce the aggregate conclusions and the saved benchmark CSVs.

\subsection{Code availability}

The implementation, experiment scripts, generated figure utilities, and reproducibility instructions will be released at \codeurl. In a final camera-ready version, this placeholder should be replaced by the public GitHub repository URL. If the submission is double-blind, the placeholder can instead be replaced by an anonymous archival link during review.
